\documentclass[10pt,a4paper]{amsart}
\usepackage[margin=19mm]{geometry}
\usepackage{amsmath,amssymb,mathtools}
\usepackage[scr=rsfs,cal=euler]{mathalfa}
\usepackage{xfrac}
\usepackage[unicode,hidelinks,bookmarks=false]{hyperref}
\newcommand{\ie}{i.e.\ }
\newcommand{\cf}{cf.\ }
\newcommand{\resp}{respectively\ }
\newcommand{\Subsection}{\subsection}
\providecommand{\nobreakdash}{\nobreak\hbox{-}}
\setlength{\emergencystretch}{2em}
\multlinegap0pt
\usepackage{bm}
\usepackage{bbm}
\usepackage{dsfont}
\usepackage{xspace}
\usepackage{cancel}
\usepackage{mathtools}
\usepackage{amsthm}
\makeatletter
\@ifundefined{theorem}{\newtheorem{theorem}{Theorem}}{}
\@ifundefined{definition}{\newtheorem{definition}[theorem]{Definition}}{}
\@ifundefined{lemma}{\newtheorem{lemma}[theorem]{Lemma}}{}
\@ifundefined{proposition}{\newtheorem{proposition}[theorem]{Proposition}}{}
\@ifundefined{remark}{\newtheorem{remark}[theorem]{Remark}}{}
\makeatother
\makeatletter
\expandafter\ifx\csname proof\endcsname\relax
\newenvironment{proof}[1][Proof]{\noindent\textbf{#1.} }{\ \rule{0.5em}{0.5em}}
\fi
\makeatother
\title[Weak almost contact structure with B-metric]{Weak almost contact structure with B-metric}
\author{Cornelia-Livia Bejan, \c{S}emsi Eken Meri\c{c}}
\date{Source: arXiv:2606.00879v1 (CC BY 4.0)}
\begin{document}
\maketitle

\noindent\emph{Source and licence.} Weak almost contact structure with B-metric. Version arXiv:2606.00879v1 (CC BY 4.0), distributed under the Creative Commons Attribution 4.0 International licence, as recorded in the arXiv licence metadata for this exact version and shown on the abstract page https://arxiv.org/abs/2606.00879v1. Full rights notice: licenses/arxiv-2606.00879-CC-BY-4.0.txt.\par\smallskip

\noindent\textit{Editorial scope.} Counted unit: the Remark whose source label is \texttt{w4} in the article (source0.tex lines 564--569, source label \texttt{w4}), delivered together with the article's own complete original proof of it (source0.tex lines 571--605, one \texttt{proof} environment of 35 lines). No mathematics is written, repaired or re-derived: statement and proof are the article's own text line for line. Required context retained uncounted: eq1 (lines 83--85); eq1a (lines 86--88); p1 (lines 102--106); eq4 (lines 233--235); p2 (lines 313--341); m2 (lines 321--323); eq101 (lines 326--329); w9 (lines 446--452); w1 (lines 537--542); w2 (lines 544--549); w5 (lines 555--557); w7 (lines 559--561). Every definition, display and label of those lines is retained unchanged. Omitted from this package and not redistributed: 1--82, 89--101, 107--232, 236--312, 324--325, 330--445, 453--536, 543--543, 550--554, 558--558, 562--563, 570--570, 606--732. Nothing inside a retained excerpt is omitted or altered: every retained body is delivered exactly as the article's lines are written, so no omission receipt of any kind accompanies this package. Citations: the retained excerpts carry no citation command at all, so no bibliography entry of the article is transcribed and no reference is redistributed. Reference adaptation: none was needed and none was made; every reference inside the delivered unit resolves inside it, so no unresolved reference or citation is produced. Printed slips kept literal: no printed word, symbol, hypothesis or number of the counted statement or of its proof was altered. Layout and class: the article's own class is \texttt{article} with packages including eurosym, amssymb, amsfonts, amsmath, mathrsfs; the wrapper is the lane's pinned \texttt{amsart} editorial wrapper at 10pt on A4, which restates the theorem-like environments the retained text needs and adds the article's own macro definitions that the retained text needs (none), with extra wrapper packages (bm, bbm, dsfont, xspace, cancel, mathtools). The delivered numbering is the unit's own. Assets: no retained span contains a figure environment, a graphics-inclusion command, a PDF-page inclusion, a table or a TikZ picture; no image file is copied into this package, the package declares no asset dependency and no image file is redistributed with it. Licence: version arXiv:2606.00879v1 of this article is distributed under the Creative Commons Attribution 4.0 International licence, as recorded in the arXiv licence metadata for this exact version and shown on the abstract page https://arxiv.org/abs/2606.00879v1. A full rights notice is delivered at licenses/arxiv-2606.00879-CC-BY-4.0.txt. Characters: every retained excerpt is pure ASCII, so the delivered engine, XeLaTeX with its default Latin Modern text font, reports no missing character.
	\begin{eqnarray}
		\varphi^2=-Q+\eta\otimes \xi, \ \label{eq1}
	\end{eqnarray}

	\begin{eqnarray}
		\eta(\xi)=1, \ \ \ \ \ Q\xi=\xi,   \label{eq1a}
	\end{eqnarray}

\begin{proposition} \label{p1}
On a weak almost contact manifold $(M^{2n+1},\varphi,Q,\xi,\eta)$, one has:\medskip

\noindent (a) $TM=\mathcal{D}\oplus Span\{\xi\}$; \ \ \ (b) $Im\varphi=\mathcal{D}$; \ \ \ \ (c) $rank \varphi=2n$; \ \ \ \ (d) $\mathcal{D}$ is $Q-$invariant, i.e., $Q(\mathcal{D})\subset\mathcal{D}$;      \ \ \ (e) $\varphi\xi=0$;   \ \ \  (f) $\eta\circ\varphi=0$;  \ \ \ (g) $\eta\circ Q=\eta$. 
\end{proposition}

	\begin{eqnarray}
		g(\varphi X,\varphi Y)=-g(X,QY)+\eta(X)\eta(Y), \ \ \ X,Y\in\chi(M).  \label{eq4}
	\end{eqnarray}

\begin{proposition} \label{p2} Let $(M^{2n+1},\varphi,Q,\xi,\eta,g)$ be a weak almost contact B-metric manifold with $\nabla$ its Levi-Civita connection. Then (i)	
		\begin{eqnarray}
		\eta &=& g(\cdot,\xi); \label{e5}
	\end{eqnarray}

\noindent (ii) $TM=\mathcal{D}\perp Span\{\xi\}$ and $g$, restricted to  both distributions is nondegenerate;\medskip

\noindent (iii)  $\nabla_X\xi$ is orthogonal to $\xi$, i.e., 
\begin{eqnarray}
	\nabla_X\xi\in\Gamma(\mathcal{D}), \ \ \ X\in\chi(M);  \label{m2}
\end{eqnarray}
	
\noindent	(iv) $\varphi$ and $Q$ are symmetric (with respect to $g$), i.e., 
	\begin{eqnarray}
		g(\varphi X,Y) = g(X,\varphi Y), \ \ \ 
			g(QX,Y) = g(X,Q Y), \ \ \ X,Y\in\chi(M); \label{eq101}
	\end{eqnarray}

\noindent (v) the (0,2)-tensor field $G$ defined as 
\begin{eqnarray}
	G(X,Y)&=&g(X,\varphi Y), \ \ \ X,Y\in\chi(M)    \label{e3}
\end{eqnarray}
is symmetric on $M$ and restricted to $\mathcal{D}$, is nondegenerate; \medskip

\noindent (vi) the (1,1)-tensor fields $\varphi$ and $Q$ commute, i.e.,
\begin{eqnarray}
	\varphi\circ Q &=& Q\circ\varphi. \label{eq9}
\end{eqnarray}
\end{proposition}

\begin{lemma}
If $(M^{2n+1},\varphi,Q,\xi,\eta, g)$ is a weak almost contact B-metric manifold, then	
\begin{eqnarray}
	g((\nabla_{\xi}\varphi)X,\varphi Y)+g((\nabla_{\xi}\varphi)Y,\varphi X)&=&-g(X,(\nabla_{\xi}Q)Y)+\eta(X)(\nabla_{\xi}\eta)(Y) \nonumber\\
	&+&\eta(Y)(\nabla_{\xi}\eta)(X), \ X,Y\in\chi(M). \label{w9}
\end{eqnarray}
\end{lemma}

\begin{definition}
A weak almost contact B-metric manifold $(M^{2n+1},\varphi,Q,\xi,\eta, g)$ is called weak nearly Sasakian B-metric manifold or weak nearly Sasaki-Norden manifold if
\begin{eqnarray}
	(\nabla_X\varphi)X&=&-g(X,X)\xi+\eta(X)X, \ \ \ X\in\chi(M).  \label{w1}
\end{eqnarray}
\end{definition}

\begin{remark}
Equivalent to (\ref{w1}), one has 	
\begin{eqnarray}
	(\nabla_X\varphi)Y+(\nabla_Y\varphi)X=-2g(X,Y)\xi+\eta(X)Y+\eta(Y)X, \ \ \ X,Y\in\chi(M).  \label{w2}
\end{eqnarray}	
\end{remark}

\begin{eqnarray}
	(\mathcal{L}_{\xi}g)(X,Y)=2g(X,\varphi Y), \ \ \ X,Y\in\chi(M),   \label{w5}
\end{eqnarray} 

\begin{eqnarray}
	(\nabla_XQ)Y=0, \ \  X\in\chi(M), \ Y\in\Gamma(\mathcal{D}). \label{w7}
\end{eqnarray} 

\begin{remark}
	In the above theorem, the statement (ii) can be written as  $\nabla_{\xi}\eta=0$ and the statement (iv) says that  the next relation is not true
	\begin{eqnarray}
		(\nabla_X\varphi)Y+(\nabla_Y\varphi)X=0, \ \ \ X,Y\in\chi(M). \label{w4}
	\end{eqnarray}
\end{remark}

\begin{proof}
(i)	We apply $\varphi$ to (\ref{w1}) and we use (\ref{eq1}), to obtain 
	\begin{eqnarray}
		\varphi\nabla_Y\varphi Y+Q\nabla_YY-\eta(\nabla_YY)\xi=\eta(Y)\varphi Y, \ \ \ Y\in\chi(M).     \nonumber
		\end{eqnarray}
	If we take $Y=\xi$, then from (\ref{m2}) it follows $Q\nabla_{\xi}\xi=0$. Hence, the nondegenerancy of $Q$ yields $\nabla_{\xi}\xi=0$. (ii) and (iii) follow from direct calculation using the fact that $\xi$ is geodesic. (iv) If we assume that (\ref{w4}) is satisfied, then from (\ref{w2}), it follows that the metric $g$ restricted to the distribution $\mathcal{D}$ is identical to zero, which contradicts the fact that $g$ is nondegenerate on the whole manifold $M$ (as we assumed $g$ is (semi-)Riemannian). We note that the relation (\ref{w2}) implies 
	\begin{eqnarray}
		(\nabla_{\xi}\varphi)X+(\nabla_X\varphi)\xi=\eta(X)\xi-X, \ \  X\in\chi(M), \label{w8}
	\end{eqnarray}
	which from Proposition \ref{p1} (e), the relation (\ref{eq4}) and (\ref{eq101}) yields
	\begin{eqnarray}
		g((\nabla_{\xi}\varphi)X,\varphi Y)=-g(\nabla_X\xi,QY)+ \eta(\nabla_X\xi)\eta(Y)-g(X,\varphi Y), \ X,Y\in\chi(M). \label{w6}
	\end{eqnarray}	
	
	Based on Proposition \ref{p2} (ii),  to obtain (v), it is enough to prove the following cases:\smallskip
	
	\noindent Case 1: If  $X=Y=\xi$, then (\ref{w5}) is satisfied, based on the relation (\ref{m2}).\smallskip
	
	\noindent Case 2: If $X\in\Gamma(\mathcal{D})$ and $Y=\xi$ (or $X=\xi$ and $Y\in\Gamma(\mathcal{D})$), then (\ref{w5}) follows from (\ref{m2}) and the parallelism of $\xi$. \smallskip
	
	\noindent Case 3: If $X,Y\in\Gamma(\mathcal{D})$, then $QY\in\Gamma(\mathcal{D})$, from Proposition \ref{p1} (d). Hence the relation (\ref{w6}) yields 
\begin{eqnarray*}
		g((\nabla_{\xi}\varphi)X,\varphi Y)=g(\xi,\nabla_XQY)-g(X,\varphi Y).
\end{eqnarray*}
A similar relation holds good by interchanging $X$ with $Y$ and if we sum these two relations, we obtain:
\begin{eqnarray*}
		&&g((\nabla_{\xi}\varphi)X,\varphi Y)+g((\nabla_{\xi}\varphi)Y,\varphi X)=g(\xi,\nabla_XQY+\nabla_YQX)-2g(X,\varphi Y)\\
		&&=g(\xi,\nabla_XY+\nabla_YX)-2g(X,\varphi Y)=(\mathcal{L}_{\xi}g)(X,Y)-2g(X,\varphi Y),
\end{eqnarray*}	
where we have used the relation	(\ref{w7}), (\ref{eq101}) and (\ref{eq1a}). On the other hand, by using (\ref{w9}), we compute the left hand side of the last relation, in which we take into account the above assertions (i) and (ii), as well as the condition (\ref{w7}), to obtain 
\begin{eqnarray*}
	g((\nabla_{\xi}\varphi)X,\varphi Y)+g((\nabla_{\xi}\varphi)Y,\varphi X)&=& 0,
\end{eqnarray*}
which gives us (\ref{w5}) and complete the proof.
\end{proof}

\end{document}
